\section{Introduction}
\label{sec:introduction-1}

Tissue organization arises from the coordinated arrangement of diverse cell types, each with distinct morphologies, phenotypes, and molecular programs~\cite{kashyap_quantification_2022}. Understanding how molecular states occur within tissues is central to many biological and clinical questions, ranging from tumor–immune interactions and stromal remodeling to spatial signaling programs and therapy resistance. 

Histopathology, based on hematoxylin and eosin (H\&E) staining of whole-slide images (WSIs), provides a rich view of tissue morphology across scales, from nuclear detail to global architectural organization. Large-scale self-supervised pretraining of vision transformers (ViTs) on millions of H\&E WSIs has led to general-purpose pathology foundation models (FMs)~\cite{chen_generalpurposefoundationmodel_2024, vorontsov_foundationmodelclinicalgrade_2024} that achieve strong performance across diverse downstream tasks, including biomarker prediction and survival modeling. However, morphology alone does not directly capture the molecular programs that determine cell identity and functional state.

Emerging spatial transcriptomics (ST) technologies promise to bridge this gap. Imaging-based ST platforms such as Xenium (10x Genomics) provide subcellular transcript coordinates for targeted gene panels within intact tissue sections~\cite{bressan_dawnspatialomics_2023}, enabling direct spatial alignment of gene expression and histology~\cite{janesick_highresolutionmapping_2023}. Public efforts such as HEST-1k~\cite{jaume_hest1kdatasetspatial_2024} and large-scale multimodal benchmarks~\cite{gindra_largescalebenchmarkcrossmodal_2025} have begun to systematize paired H\&E and ST datasets across tissues and panels. In this study, we therefore complement our internal cohort with organ-specific HEST-1k-derived cohorts and HESCAPE benchmark subsets~\cite{jaume_hest1kdatasetspatial_2024,gindra_largescalebenchmarkcrossmodal_2025}, and further consider BEAT as an additional setting for panel-restricted and highly variable gene prediction. Together, these developments create the opportunity to learn unified morpho-molecular representations of tissue ecosystems.

\subsection*{Multimodal Learning in Spatial Omics}
\label{sec:multimodal-learning-in-spatial-omics}

Recent work has explored multiple paradigms for integrating morphology and transcriptomics. A first line of research focuses on predicting gene expression directly from H\&E images (e.g., ST-Net~\cite{he_integratingspatialgene_2020}, HisToGene~\cite{pang_leveraginginformationspatial_2021}, DeepSpot~\cite{nonchev_deepspotleveragingspatial_2025}), treating transcriptomics as a supervised regression target. More recent approaches adopt CLIP-style contrastive pretraining~\cite{radford_learningtransferablevisual_2021}, aligning image and gene-expression encoders in a shared latent space (e.g., OmiCLIP~\cite{chen_visualomicsfoundation_2025}, PathOmCLIP~\cite{lee_pathomclipconnectingtumor_2024}, \linebreak BLEEP~\cite{xie_spatiallyresolvedgene_2023}).


While contrastive alignment improves cross-modal retrieval, it does not constitute true multimodal fusion. Transcriptomic measurements act as supervisory signals that shape image encoders during pretraining, but the resulting representations remain unimodal at inference. Cross-modal interactions are restricted to pooled embedding spaces rather than being modeled jointly at the level of spatially resolved tokens. 

Large-scale benchmarks highlight limitations of this paradigm. In HESCAPE~\cite{gindra_largescalebenchmarkcrossmodal_2025}, contrastive pretraining improves mutation classification but degrades direct gene expression prediction compared to strong baseline encoders. These results suggest that alignment objectives may fail to leverage molecular information beyond what is already encoded in morphology, potentially due to batch effects and panel heterogeneity. Importantly, most evaluated architectures rely on late-fusion or embedding-level alignment, leaving spatially resolved token-level interaction completely unexplored.

\subsection*{Limitations of Cell-Centric Reduction}
\label{sec:limitations-of-cell-centric-reduction}

A second limitation in current spatial omics modeling lies in the standard analysis workflow itself. Imaging-based ST data are typically processed in a cell-centric manner: transcripts are assigned to segmentation masks, aggregated into per-cell expression matrices, clustered, and annotated into cell types. While powerful, this pipeline reduces rich spatial information (morphology, shape, neighborhood structure, and subcellular transcript distribution) to a single expression vector per segmented cell. 

Whole-cell segmentation is time-consuming and error-prone, particularly in densely packed regions or for morphologically complex cell types such as fibroblasts. Transcript diffusion and partial segmentation further complicate quantification~\cite{bilous_transcriptscellsdissecting_2025}. Moreover, recent studies indicate that cellular phenotypes are often spatially modulated, with location-dependent subtypes emerging within what were previously considered homogeneous populations. Reducing spatial omics images to expression matrices therefore risks obscuring morpho-spatial determinants of cellular identity.

To fully exploit spatial transcriptomics, we require architectures that jointly model morphology, molecular profiles, and spatial context without relying exclusively on segmentation-derived abstractions.

\subsection*{Two Tales: From Gene Prediction to Cellular Phenotyping}
\label{sec:from-gene-prediction-to-phenotyping}

We evaluate our multimodal framework on two complementary tasks that reflect distinct biological questions and difficulty levels, each motivated by practical realities of the Xenium platform.

\paragraph{Task 1: Add-on gene prediction.}
Xenium panels are structured around two complementary components: a \emph{core panel}, supplied by the platform provider and covering well-established gene signatures for common cell types and pathways, and a user-configurable \emph{add-on panel} of study-specific probes selected by researchers based on their biological question. This two-tier panel architecture arises naturally from the platform design and has direct practical consequences: researchers must decide upfront which add-on genes to include, often without knowing which of those can already be reliably inferred from core measurements or from tissue morphology.

Predicting add-on gene expression from the core panel and histology is therefore both technically meaningful and practically motivated. First, it enables \emph{panel optimization}: if a given add-on gene can be accurately predicted from the core panel alone, it may be expendable in future experiments. Second, it allows \emph{panel augmentation}: genes not originally measured may be recoverable in silico from existing data. We implement this task by predicting the log$_{1p}$-transformed expression of $G_{\text{add}} = 100$ add-on genes from the $G_{\text{core}} = 380$ gene core panel and H\&E histology.

\paragraph{Task 2: Cell type composition prediction.}
Cell type composition represents a higher-level phenotype that encodes the collective identity of cells within a tissue region. Unlike individual gene expression values, cell type abundance integrates multiple coordinated gene programs that collectively define a cell identity. This integration happens at the level of individual cells and is therefore strongly modulated by local co-expression patterns that may be averaged out or obscured when gene expression is summarized at the tile level.

This should make cell type prediction inherently harder than tile-level gene prediction: a model must not merely recover aggregate molecular profiles but must respect the local co-expression structure that distinguishes, for example, a cytotoxic T cell from an NK cell or a macrophage from a monocyte. At the same time, this task carries direct biological and clinical value. Cell type annotations are labor-intensive, requiring whole-slide segmentation, transcript assignment, clustering, and expert curation. A model that can predict cell type abundance from partial transcriptomic input and H\&E images, even at a coarse tile level, would substantially reduce this burden and extend annotation capacity to datasets where only histology or partial panel measurements are available. We note that our current formulation provides tile-level predictions, and we are actively working toward more granular, cell-level predictions (discussed in Section~\ref{sec:limitations-and-future-directions}).

We therefore introduce cell type count prediction as a complementary and more stringent evaluation task. Specifically, we predict log$_{1p}$-transformed counts of 39 phenotypically defined cell types within each tile of our samples. Unlike gene imputation, this task cannot be trivially resolved through single-gene correlations; it requires the model to capture coordinated morpho-molecular patterns that collectively define cellular ecosystems.

\subsection*{Our Contribution}
\label{sec:our-contribution}

In this work, we introduce the first token-level early-fusion multimodal transformer architecture for joint representation learning of histology and Xenium spatial transcriptomics. Our approach directly exploits subcellular transcript readouts by injecting transcript-derived tokens into the ViT token sequence prior to any shared transformer processing. This enables all transformer layers to operate on jointly contextualized morphology--transcript representations, rather than restricting cross-modal interaction to aggregated feature spaces.

We evaluate our architecture on the two tasks described above: (i) add-on gene prediction from the 380-gene core panel, and (ii) cell type count prediction for 39 phenotypically defined cell populations, evaluated separately for core, add-on, and full panel inputs. While multimodal fusion yields consistent improvements in gene prediction over unimodal baselines, its benefit is most pronounced in the phenotyping task, where token-level early fusion substantially outperforms unimodal and late-fusion baselines. We further show that increasing model capacity (through a larger ViT backbone and deeper expression encoder) improves fusion performance but provides negligible benefit to morphology-only models, suggesting that the additional expressiveness is specifically needed for modeling cross-modal interactions. Finally, per-cell-type analysis reveals that unimodal models exhibit characteristic failure modes on specific cell populations, motivating the need for joint multimodal representations.

These results demonstrate that spatially aligned, multimodal interaction is critical when modeling higher-order tissue phenotypes. Beyond alignment, our architecture provides a principled framework for learning unified morpho-molecular representations that reflect cellular composition and spatial organization, moving closer toward integrative foundation models for spatial biology.
